\section{The Input Block for Surfaces}
\label{sec2.3}
%\addcontentsline{toc}{section}{\protect\numberline{2.3}{The Input
%Block for Surfaces}}
Grid surfaces may be assigned special properties (e.g.: reflecting,
absorbing, periodicity, modified cell index switching, etc.).
Their definition is described in \cref{sec2.3a}.

In addition to these grid surfaces also additional surfaces can be
seen by the histories (vacuum boundaries, special surfaces for
scoring fluxes, diagnostic surfaces, \dots).
Such surfaces can be general linear or second order surfaces in 3D space.
Their definition and properties are described in \cref{sec2.3b}.

\subsection{The Input Block for ``Non-default Standard Surfaces"}
\label{sec2.3a}


\medskip

\textbf{General remarks}

\medskip

Some of the ``standard grid surfaces'' RSURF, PSURF and TSURF may be
defined as reflecting or absorbing (rather than the default
transparent character of co-ordinate surfaces), or as transparent
but with other switching features (cell indexing, transition into
additional cell region etc.) than the default settings. E.g.,
surface averaged tallies (see block 11 (\cref{sec2.11}), end of 
\cref{sec3.2} and \cref{tab3}) are only evaluated on either
``Additional Surfaces'' (\cref{sec2.3b}) or on these ``Non
default Standard Surfaces'', whereas they are not computed for those
standard grid surfaces which are not explicitly identified in this
block.

The geometrical properties of these ``Non default Standard Surfaces" are
specified here. This is described below.

\textbf{Particle-surface interaction models}

Although the (local) data for particle-surface interaction models
(PSI-models) for each specific surface can be read in this input
block, their meaning is described in block 6 (\cref{sec2.6}) together
with the globally valid input data for particle-surface
interactions.

In any particular run there are NLIM ``additional surfaces", see
\cref{sec2.3b}. For many arrays containing surface
properties the ``non-default standard surfaces" are stored after the
additional surfaces. E.g.\ FLAG(NLIM+4) would store the data FLAG for
non-default standard surface no. 4.

\medskip

\textsl{*** 3A. Data for ``Non-default Standard Surfaces''}
\begin{verbatim}
 NSTSI
 DO 31 ISTSI=1,NSTSI
\end{verbatim}
 unless otherwise stated, N = NLIM+ISTSI is the index of the following
 surface data for surface ISTSI:
\begin{verbatim}
    TXTSFL
    ISTS, IDIMP, INUMP(ISTSI,IDIMP),
   .IRPTA1, IRPTE1, IRPTA2, IRPTE2, IRPTA3, IRPTE3
    ILIIN, ILSIDE, ILSWCH, ILEQUI, ILTOR, ILCOL,
   .ILFIT ILCELL ILBOX ILPLG
\end{verbatim}
optional for non transparent surfaces (ILIIN $>$ 0):

The next 3 (or 4) lines comprise the block for local particle-surface
interaction data
(in analytic terminology: the boundary condition at this surface element).
If they are omitted, the default particle-surface model is activated for this
particular surface element, see \cref{sec2.6}.
\begin{verbatim}
    ILREF  ILSPT ISRS  ISRC
    ZNML EWALL EWBIN TRANSP(1,N) TRANSP(2,N) FSHEAT
    RECYCF  RECYCT  RECPRM  EXPPL  EXPEL  EXPIL
    RECYCS  RECYCC  SPTPRM  (this line may be omitted,
        then: default sputter model, see \cref{sec2.6})
 31 CONTINUE
\end{verbatim}


also, optional for non-transparent surfaces (ILIIN $>$ 0), and
alternative to the local surface interaction data block mentioned
above:

read a surface interaction model identifier to link one of the surface
``local reflection
models'' defined in block 6 (\cref{sec2.6}) below to this current surface.
The presence of such a link is identified by the string $'SURFMOD_{ }'$,
followed
by a name $'modname'$. Later, in input block 6 (\cref{sec2.6}),
a ``local reflection model'' block with that name
$'modname'$ must be included. This option allows
quick changes of particle-surface interaction parameters, affecting
many surfaces at once,
by changes in just one location of the input file.

\begin{verbatim}
   SURFMOD_modname
\end{verbatim}
\medskip

\textbf{Meaning of the Input Variables for ``non-default standard
surfaces"}

\medskip

\begin{list}{ }{}
\item[\textbf{NSTSI }]
Total number of non-default standard surfaces that do not act as
prescribed by the default transparent standard co-ordinate surface model.
\item[\textbf{TXTSFL }]
Text to characterize a surface (``name of the surface'') on the
printout file.
\item[\textbf{ISTS }]
irrelevant; labelling index
\item[\textbf{IDIMP }]
flag to identify mesh from which this particular surface is
chosen.
\begin{list}{ }{}
\item[{= 1 }]
surface from the x- (radial) standard mesh (RSURF)
Note that for the unstructured grid options NLTRI and NLTET (i.e.\ for the 2D triangular grid option
or for the general 3D grids of tetrahedra) all surfaces are referred to as 1st grid (x or radial) surfaces,
by abuse of language.
\item[{= 2 }]
surface from the y- (poloidal) standard mesh (PSURF)
\item[{= 3 }]
surface from the z- (toroidal) standard mesh (TSURF)
\end{list}
\item[\textbf{INUMP }]
Number of the surface in mesh RSURF, PSURF or TSURF respectively
\item[\textbf{IRPTA,IRPTE}] ~

Only a subregion of the surface acts by the ``non-default options"
specified for this particular surface. This subregion is defined
by these flags.

If JMP is a surface from the first mesh, then IRPTA2$\rightarrow
$IRPTE2 and
IRPTA3$\rightarrow
$IRPTE3 are the surface index ranges of the 2nd and 3rd mesh,
respectively, for which this surface acts as non-default surface.
IRPTA1 and IRPTE1 are irrelevant.

If JMP is a surface from the 2nd mesh, then IRPTA1$\rightarrow
$IRPTE1 and
IRPTA3$\rightarrow
$IRPTE3 are the surface index ranges of the 1st and 3rd mesh,
respectively, for which this surface acts as non-default surface.
IRPTA2 and IRPTE2 are irrelevant.

If JMP is a surface from the 3rd mesh, then IRPTA1$\rightarrow$
IRPTE1 and
IRPTA2$\rightarrow$
IRPTE2 are the surface index ranges of the 1st and 2rd mesh,
respectively, for which this surface acts as non-default surface.
IRPTA3 and IRPTE3 are irrelevant.
\end{list}

The third card
ILIIN, \dots, ILPLG is identical to the corresponding card for ``additional surfaces",
 see block 3b (\cref{sec2.3b}) below,
``Input Data for Additional Surfaces''. One exception is the flag ILSIDE, which controls
the sign of the surface normal vector (hence the orientation of a surface). In case of unstructured
``standard grids" NLTRI or NLTET (triangles in 2D and tetrahedrons in 3D) there is no well defined
default surface orientation and the flag ILSIDE is irrelevant in such cases.

\medskip

The input data in the block for the local reflection model are
described below, see \cref{sec2.6} : ``Input Data for Particle-Surface Interaction
Models''.
